\paragraph{R4-O3 (Validation may be expensive under complex constraints; discuss validation cost).}

\begin{itemize}[leftmargin=1.2em, itemsep=1pt, topsep=2pt]
\item \textit{O3: ``\dots validation itself may be expensive under complex constraints. The authors are encouraged to further discuss the validation cost.''}
\end{itemize}
\smallskip

\noindent\textbf{Response.}
The coarse-to-fine design (Section~2.2) bounds validation cost. Cheap attribute/metric checks are pushed into candidate generation (Step~2.a), and expensive execution-level (temporal, graph) or semantic (text, image) validation is applied only to the candidate pool $\mathcal{A}_\omega(q)$ of size $m\ge k$; validation cost therefore scales with $|\mathcal{A}_\omega|$, not $|\mathcal{D}|$. Because the pool is returned \emph{``in nondecreasing cost order for early pruning,''} validation can stop once $k$ admissible answers are confirmed. Table~\ref{tab:scale} shows the online selection cost is near-constant in $N$, consistent with validation being confined to the small pool.

\noindent\emph{Revision to Section~3 (validation cost):}
\begin{quote}
\textcolor{blue}{The coarse-to-fine design bounds the cost of constraint enforcement and feasibility validation. Cheap attribute- and metric-level checks are performed during candidate generation (Step~2.a), pruning most infeasible instances before expensive validation. Execution-level (temporal, graph) and semantic (text, image) validation is then applied only to the candidate pool $\mathcal A_\omega(q)$, so validation scales with $|\mathcal A_\omega|$ rather than $|\mathcal D|$. Because candidates are processed in nondecreasing cost order, validation terminates once $k$ admissible answers are found. Increasing $m$ improves the likelihood of finding $k$ feasible counterfactuals at the cost of additional validation, providing a controllable accuracy--latency trade-off (Principle~P3 in Section~3.2).}
\end{quote}
